% GENERATED FILE. Edit canonical lessons, then run npm run generate:books.
% canonical-source-hash: fnv1a64:2296e9611a650e11
% canonical-lessons: BN-C398-srot, BN-C398-soikot, BN-C398-jeti, BN-C398-dub-deoya, BN-C398-chhitano

\chapter{Currents and Beaches}
\label{ch:bn-currents-and-beaches}
\hlchaptermodality{398}

\begin{chapteropening}
\textbf{By the end of this chapter:} \emph{I can talk about currents, beaches and jetties, and a dip in the water.}

Complete the last lesson of chapter 398: I can talk about currents, beaches and jetties, and a dip in the water.
\end{chapteropening}

\section[\texorpdfstring{\bn{স্রোত} (srot)}{স্রোত (srot)}]{\bn{স্রোত} (srot) — a current; a stream}

\label{lesson:BN-C398-srot}



\begin{quote}
Before the new one: say the Bengali for a bank (of a pond or river), then the Bengali for to moor.
\end{quote}

\subsection*{You'll want to know: \bn{স্রোত}}
\textbf{\bn{স্রোত}} — \emph{srot} — \textquotedblleft{}a current; a stream\textquotedblright{}.

\begin{quote}
\bn{নদীর} \bn{স্রোতে} \bn{নৌকাটা} \bn{দ্রুত} \bn{চলল।}
\end{quote}

\begin{quote}
\emph{nôdīr srote noukāṭā drutô chôllô.}
\end{quote}

\begin{quote}
The boat moved fast on the river's current.
\end{quote}

Also a flow of people or things: \bn{মানুষের} \bn{স্রোত}, a stream of people.

\begin{cousinweb}
\textbf{\bn{স্রোত}} (\emph{srot}) is Sanskrit \emph{srotas}, \textquotedblleft{}the current or bed of a river, a stream\textquotedblright{} (Monier-Williams), from \emph{sru}, to flow.
\end{cousinweb}

\begin{grammarlens}[title={What you've built}]
One more word for getting around.
\end{grammarlens}

\subsection*{Your turn}
\begin{itemize}
\raggedright
  \item \emph{Say it:} \emph{srot}
  \item \emph{Say it:} \emph{srot}, once more
  \item \emph{Say it:} \emph{nôdīr srote noukāṭā drutô chôllô.} — The boat moved fast on the river's current.
  \item \emph{Recall:} say the Bengali for a bank (of a pond or river), then the Bengali for to moor, then say \emph{srot} again
\end{itemize}

\begin{tcolorbox}[breakable,colback=teal!4,colframe=teal!35!black,title={Before you move on}]
What does \bn{স্রোত} mean? (\textquotedblleft{}A current; a stream\textquotedblright{}.) Say it once more.
\end{tcolorbox}
\section[\texorpdfstring{\bn{সৈকত} (soikôt)}{সৈকত (soikôt)}]{\bn{সৈকত} (soikôt) — a beach}

\label{lesson:BN-C398-soikot}



\begin{quote}
Before the new one: say the Bengali for to moor, then the Bengali for a current.
\end{quote}

\subsection*{You'll want to know: \bn{সৈকত}}
\textbf{\bn{সৈকত}} — \emph{soikôt} — \textquotedblleft{}a beach\textquotedblright{}.

\begin{quote}
\bn{সন্ধ্যায়} \bn{আমরা} \bn{সৈকতে} \bn{হাঁটলাম।}
\end{quote}

\begin{quote}
\emph{sôndhyāy āmrā soikôte hā̃ṭlām.}
\end{quote}

\begin{quote}
In the evening we walked on the beach.
\end{quote}

The first vowel sign is \bn{ৈ}, the vowel oi; in talk people also say \bn{সমুদ্রের} \bn{ধার}.

\begin{cousinweb}
\textbf{\bn{সৈকত}} (\emph{soikôt}) is Sanskrit \emph{saikata}, \textquotedblleft{}sandy, made of sand\textquotedblright{} (Monier-Williams), from \emph{sikatā}, sand: a beach is the sandy place.
\end{cousinweb}

\begin{grammarlens}[title={What you've built}]
One more word for getting around.
\end{grammarlens}

\subsection*{Your turn}
\begin{itemize}
\raggedright
  \item \emph{Say it:} \emph{soikôt}
  \item \emph{Say it:} \emph{soikôt}, once more
  \item \emph{Say it:} \emph{sôndhyāy āmrā soikôte hā̃ṭlām.} — In the evening we walked on the beach.
  \item \emph{Recall:} say the Bengali for to moor, then the Bengali for a current, then say \emph{soikôt} again
\end{itemize}

\begin{tcolorbox}[breakable,colback=teal!4,colframe=teal!35!black,title={Before you move on}]
What does \bn{সৈকত} mean? (\textquotedblleft{}A beach\textquotedblright{}.) Say it once more.
\end{tcolorbox}
\section[\texorpdfstring{\bn{জেটি} (jeṭi)}{জেটি (jeṭi)}]{\bn{জেটি} (jeṭi) — a jetty}

\label{lesson:BN-C398-jeti}



\begin{quote}
Before the new one: say the Bengali for a current, then the Bengali for a beach.
\end{quote}

\subsection*{You'll want to know: \bn{জেটি}}
\textbf{\bn{জেটি}} — \emph{jeṭi} — \textquotedblleft{}a jetty\textquotedblright{}.

\begin{quote}
\bn{জেটিতে} \bn{দাঁড়িয়ে} \bn{আমরা} \bn{জাহাজ} \bn{দেখলাম।}
\end{quote}

\begin{quote}
\emph{jeṭite dā̃ṛiye āmrā jāhāj dekhlām.}
\end{quote}

\begin{quote}
We stood on the jetty and watched the ships.
\end{quote}

\begin{cousinweb}
\textbf{\bn{জেটি}} (\emph{jeṭi}) is borrowed from English \emph{jetty}, a landing stage built out into the water.
\end{cousinweb}

\begin{grammarlens}[title={What you've built}]
One more word for getting around.
\end{grammarlens}

\subsection*{Your turn}
\begin{itemize}
\raggedright
  \item \emph{Say it:} \emph{jeṭi}
  \item \emph{Say it:} \emph{jeṭi}, once more
  \item \emph{Say it:} \emph{jeṭite dā̃ṛiye āmrā jāhāj dekhlām.} — We stood on the jetty and watched the ships.
  \item \emph{Recall:} say the Bengali for a current, then the Bengali for a beach, then say \emph{jeṭi} again
\end{itemize}

\begin{tcolorbox}[breakable,colback=teal!4,colframe=teal!35!black,title={Before you move on}]
What does \bn{জেটি} mean? (\textquotedblleft{}A jetty\textquotedblright{}.) Say it once more.
\end{tcolorbox}
\section[\texorpdfstring{\bn{ডুব} \bn{দেওয়া} (ḍub deoyā)}{ডুব দেওয়া (ḍub deoyā)}]{\bn{ডুব} \bn{দেওয়া} (ḍub deoyā) — to take a dip; to dive under}

\label{lesson:BN-C398-dub-deoya}



\begin{quote}
Before the new one: say the Bengali for a beach, then the Bengali for a jetty.
\end{quote}

\subsection*{You'll want to know: \bn{ডুব} \bn{দেওয়া}}
\textbf{\bn{ডুব} \bn{দেওয়া}} — \emph{ḍub deoyā} — \textquotedblleft{}to take a dip; to dive under\textquotedblright{}.

\begin{quote}
\bn{গরমে} \bn{অরুণ} \bn{পুকুরে} \bn{ডুব} \bn{দিল।}
\end{quote}

\begin{quote}
\emph{gôrôme ôrun pukure ḍub dilô.}
\end{quote}

\begin{quote}
In the heat Arun took a dip in the pond.
\end{quote}

Past (he/she …): \bn{তিনি} \bn{পুকুরে} \bn{ডুব} \bn{দিলেন}, \emph{tini pukure ḍub dilen}.

\bn{ডুব} is a plunge under water; \bn{ডোবা}, to sink, is the same root.

\begin{cousinweb}
\textbf{\bn{ডুব} \bn{দেওয়া}} (\emph{ḍub deoyā}) joins \bn{ডুব}, \textquotedblleft{}a plunge under water\textquotedblright{}, from the reconstructed \emph{ḍubb}, \textquotedblleft{}sink\textquotedblright{}, under which CDIAL lists Bengali \emph{ḍubā}, \textquotedblleft{}to sink\textquotedblright{} (CDIAL 5561), to \bn{দেওয়া}, to give.
\end{cousinweb}

\begin{grammarlens}[title={What you've built}]
One more word for getting around.
\end{grammarlens}

\subsection*{Your turn}
\begin{itemize}
\raggedright
  \item \emph{Say it:} \emph{ḍub deoyā}
  \item \emph{Say it:} \emph{ḍub deoyā}, once more
  \item \emph{Say it:} \emph{gôrôme ôrun pukure ḍub dilô.} — In the heat Arun took a dip in the pond.
  \item \emph{Recall:} say the Bengali for a beach, then the Bengali for a jetty, then say \emph{ḍub deoyā} again
\end{itemize}

\begin{tcolorbox}[breakable,colback=teal!4,colframe=teal!35!black,title={Before you move on}]
What does \bn{ডুব} \bn{দেওয়া} mean? (\textquotedblleft{}To take a dip; to dive under\textquotedblright{}.) Say it once more.
\end{tcolorbox}
\section[\texorpdfstring{\bn{ছিটানো} (chhiṭāno)}{ছিটানো (chhiṭāno)}]{\bn{ছিটানো} (chhiṭāno) — to splash; to sprinkle}

\label{lesson:BN-C398-chhitano}



\begin{quote}
Before the new one: say the Bengali for a jetty, then the Bengali for to take a dip.
\end{quote}

\subsection*{You'll want to know: \bn{ছিটানো}}
\textbf{\bn{ছিটানো}} — \emph{chhiṭāno} — \textquotedblleft{}to splash; to sprinkle\textquotedblright{}.

\begin{quote}
\bn{শিশুরা} \bn{সৈকতে} \bn{জল} \bn{ছিটাল।}
\end{quote}

\begin{quote}
\emph{shishurā soikôte jôl chhiṭālô.}
\end{quote}

\begin{quote}
The children splashed water on the beach.
\end{quote}

Past (he/she …): \bn{তিনি} \bn{জল} \bn{ছিটালেন}, \emph{tini jôl chhiṭālen}.

\begin{cousinweb}
\textbf{\bn{ছিটানো}} (\emph{chhiṭāno}) is inherited: CDIAL lists Bengali \emph{chiṭā}, \textquotedblleft{}to scatter, to sprinkle\textquotedblright{}, under the reconstructed \emph{chiṭ}, \textquotedblleft{}a sudden movement, a splash\textquotedblright{} (CDIAL 5035).
\end{cousinweb}

\begin{grammarlens}[title={What you've built}]
One more word for getting around.
\end{grammarlens}

\subsection*{Your turn}
\begin{itemize}
\raggedright
  \item \emph{Say it:} \emph{chhiṭāno}
  \item \emph{Say it:} \emph{chhiṭāno}, once more
  \item \emph{Say it:} \emph{shishurā soikôte jôl chhiṭālô.} — The children splashed water on the beach.
  \item \emph{Recall:} say the Bengali for a jetty, then the Bengali for to take a dip, then say \emph{chhiṭāno} again
\end{itemize}

\begin{tcolorbox}[breakable,colback=teal!4,colframe=teal!35!black,title={Before you move on}]
What does \bn{ছিটানো} mean? (\textquotedblleft{}To splash; to sprinkle\textquotedblright{}.) Say it once more.
\end{tcolorbox}
